\documentclass[11pt]{article}
\usepackage{fontspec}
\usepackage[french,english]{babel}
\usepackage{amsmath}
\usepackage{amssymb}
\usepackage{amsthm}
\usepackage{mathrsfs}
\usepackage{tikz-cd}
\usepackage{enumitem}
\usetikzlibrary{babel}
\newcommand{\printedpage}[1]{\marginpar{\footnotesize p.~#1}}
\newcommand{\msp}[1]{\marginpar{\footnotesize [#1]}}
\newcommand{\sX}{\mathcal{X}} % resolved against PDF: CMSY calligraphic X
\begin{document}

\printedpage{9}
% head: A. Grothendieck : Dérivateurs, Ch. II, Édition des Universités de Montpellier II et Paris VII
% foot: Version 8 juillet 2001

\msp{page 14}

In summary, we have the

\noindent\textbf{Proposition 2.1.} \textit{Let $X$ be an object of} Cat. \textit{Equivalent conditions:}
\begin{enumerate}[label={},leftmargin=3em]
\item[\textit{a)}] \textit{The associated topos} $\operatorname{Top}(X)$ \textit{is locally irreducible, i.e. (being locally connected) for every object $F$ of $X^{\wedge}$ and every subobject $F'$ of $F$, if $F$ is connected, so is $F'$.}
\item[\textit{a$'$)}] \textit{For every $X'$ étale over $X$, if $X'$ is connected, every open of $X'$ is connected.}
\item[\textit{a$''$)}] \textit{Every open of a localized category $X/a$ ($a \in \operatorname{Ob} X$) is connected (i.e. $X/a$ is an irreducible category, i.e. the topos} $\operatorname{Top}(X/a)$ \textit{is).}
\item[\textit{b)}] \textit{For every object $a$ of $X$,} $\operatorname{Esp}(X/a)$ \textit{is an irreducible topological space.}
\item[\textit{c)}] \textit{The categories $X/a$ are weakly left filtering (or weakly cofiltering), i.e. for two objects $\xi$, $\eta$ of such a category,} \texttt{[there]} \textit{exists an object $\zeta$ and two arrows}
\begin{tikzcd}[ampersand replacement=\&, row sep=tiny, column sep=small, baseline=(\tikzcdmatrixname.center)]
 \& \xi \\
\zeta \arrow[ur] \arrow[dr] \& \\
 \& \eta
\end{tikzcd}
\textit{(i.e. the ordered set Ord($X/a$) associated with $X/a$ is left filtering)}
\item[\textit{c$'$)}] \textit{For every diagram in solid lines in $X$,}
\begin{tikzcd}[ampersand replacement=\&, row sep=tiny, column sep=small, baseline=(\tikzcdmatrixname.center)]
 \& x \arrow[dr] \& \\
z \arrow[ur, dotted] \arrow[dr, dotted] \& \& s \\
 \& y \arrow[ur] \&
\end{tikzcd}
\textit{, we can complete it to a} commutative \textit{square as indicated. (Condition} PS 1 \textit{for locally filtering categories.)}
\end{enumerate}

\noindent These conditions imply that every connected component $X_{\alpha}$ of $X$ is irreducible, or what amounts to the same thing, that $\operatorname{Esp}(X_{\alpha})$ is irreducible,

\msp{page 15}

\noindent but are not implied by the latter.

Counterexample, the 0-connected category $0 \underset{\beta}{\overset{\alpha}{\rightrightarrows}} 1$, with $\alpha \neq \beta$. We have $X/1 \simeq$
\begin{tikzcd}[ampersand replacement=\&, row sep=tiny, column sep=small, baseline=(\tikzcdmatrixname.center)]
0 \arrow[dr] \& \\
 \& 1 \\
0' \arrow[ur] \&
\end{tikzcd}, which is \textit{not} irreducible, since $\Phi \backslash \{1\} \simeq \varepsilon$ is a disconnected open \texttt{[}$\Phi = X/1$\texttt{,} $\varepsilon = \bullet\ \bullet$\texttt{]}.

I would like a ``topossic'' condition on $X$, equivalent to $X$ ``left pseudo-filtering'', i.e. the connected components filtering, i.e. PS 1\textsuperscript{opp} (condition c$'$) above) + PS 2\textsuperscript{opp} (double arrow condition), and which moreover implies that every 0-connected object of $X^{\wedge}$ (i.e. every object $X'$ of Cat étale over $X$ and 0-connected) is $W_{\infty}$-aspheric. And I am thinking of the notion of ``total asphericity'' of Pursuing Stacks.

\noindent\textbf{Definition 2.2.} (\footnote[13]{Cf. summary table pages 97, 98 and especially 99.})\ \ Let $X$ be a category.
\begin{enumerate}[label={},leftmargin=3em]
\item[a)] We say that $X$ is \textit{weakly filtering} (\footnote[14]{rather say \textit{pseudo-filtering}, reserve ``weakly filtering'' for ``locally weakly filtering'' (or \textit{irreducible}) + connected, condition PS $1'_{\mathrm{g}}$.}) if $X \neq \emptyset$, and if for $a, b \in \operatorname{Ob} X$, $\exists\, c$ and arrows $c$
\begin{tikzcd}[ampersand replacement=\&, row sep=tiny, column sep=small, baseline=(\tikzcdmatrixname.center)]
 \& a \\
{} \arrow[ur, phantom] \& \\
 \& b
\end{tikzcd}
\hspace{-2.6em}% layout aid only
\begin{tikzcd}[ampersand replacement=\&, row sep=tiny, column sep=small, baseline=(\tikzcdmatrixname.center)]
 \& {} \\
{} \arrow[ur] \arrow[dr] \& \\
 \& {}
\end{tikzcd}
, i.e. $X/[a] \cap X/[b] \neq \emptyset$. (\textbf{NB} $X$ is weakly left filtering (\footnote[15]{pseudo-filtering.}) if and only if $\operatorname{Ord}(X)$ is, i.e. is left filtering. \textbf{ReNB} $X$ weakly filtering $\Longrightarrow$ $X$ 0-connected.)

\item[b)] \printedpage{10}
% head: A. Grothendieck : Dérivateurs, Ch. II, Édition des Universités de Montpellier II et Paris VII
% foot: Version 8 juillet 2001
We say that $X$ is \textit{locally weakly left filtering} or more simply left prefiltering if the categories $X/x$ ($x \in \operatorname{Ob} X$) are weakly left filtering, i.e. if $X$ satisfies the condition: $\forall$ diagrams
\begin{tikzcd}[ampersand replacement=\&, row sep=tiny, column sep=small, baseline=(\tikzcdmatrixname.center)]
a \arrow[dr] \& \\
 \& x \\
b \arrow[ur] \&
\end{tikzcd}, $\exists\, c$ in $X$ and arrows from $c$ to $a$, $b$ making the square \textit{commutative}
\[
\begin{tikzcd}[ampersand replacement=\&, row sep=tiny, column sep=small]
 \& a \arrow[dr] \& \\
c \arrow[ur, dotted] \arrow[dr, dotted] \& \& x \\
 \& b \arrow[ur] \&
\end{tikzcd}
\]
(\footnote[16]{Condition PS $1_{\mathrm{g}}$.}). (\textbf{NB} This condition implies that $\operatorname{Ord}(X)$ is locally weakly filtering, but is not equivalent to the latter. Example, $X = 0 \underset{\beta}{\overset{\alpha}{\rightrightarrows}} 1$, $X/1 =$
\begin{tikzcd}[ampersand replacement=\&, row sep=tiny, column sep=small, baseline=(\tikzcdmatrixname.center)]
0' \arrow[dr] \& \\
 \& 1 \\
0'' \arrow[ur] \&
\end{tikzcd}
is not weakly left filtering.)

\msp{page 16}

\textbf{NB} If $X$ locally weakly left filtering, every $X'$ étale over $X$ likewise.

\textbf{NB} (\footnote[17]{\textbf{False.}}) $X$ locally weakly left filtering $\Longleftrightarrow$ the connected components are weakly left filtering, or again if $X$ is 0-connected, $X$ locally weakly left filtering $\Longleftrightarrow$ $X$ is weakly left filtering (\footnote[18]{pseudo-filtering.}). We have indeed $\Longleftarrow$ tautologically (\footnote[19]{\textbf{false,} thus $0 \underset{\beta}{\overset{\alpha}{\rightrightarrows}} 1$ is 0-connected and weakly left filtering, but not locally weakly left filtering because
\begin{tikzcd}[ampersand replacement=\&, row sep=tiny, column sep=small, baseline=(\tikzcdmatrixname.center)]
0 \arrow[dr, "\alpha"] \& \\
 \& 1 \\
0 \arrow[ur, "\beta"'] \&
\end{tikzcd}
cannot be completed to a \textit{commutative} square.}), and $\Longrightarrow$ by showing that if PS 1 \texttt{[is]} satisfied, then $PS(a,b)$ defined by $X/[a] \cap X/[b] \neq \emptyset$ is an equivalence relation, implied by the relation $C(a,b)$: $\{\operatorname{Hom}(a,b) \neq \emptyset$ \textit{or} $\operatorname{Hom}(b,a) \neq \emptyset\}$, hence the equivalence relation (homotopy) generated by $C(a,b)$ is finer than $PS(a,b)$.

\item[c)] We say that $X$ is \textit{left filtering}, or \textit{cofiltering,} if it is weakly left filtering (\footnote[20]{pseudo-filtering.}) (hence \textit{nonempty}), and satisfies (in addition to condition PS $1'_{\mathrm{g}}$ of a)) the following condition PS $2_{\mathrm{g}}$:

For every double arrow $a \underset{v}{\overset{u}{\rightrightarrows}} b$, \texttt{[there]} exists \texttt{[an]} arrow $c \xrightarrow{w} a$ such that $uw = vw$ (\footnote[21]{i.e. three conditions a) $X \neq \emptyset$ b) usual PS $1'$ c) PS 2 (double arrow).}).

It amounts to the same to say that $X$ satisfies the three conditions a$'$) $X$ is 0-connected (stronger than $X \neq \emptyset$), b$'$) PS 1, i.e. $X$ locally weakly left filtering and c$'$) PS 2 as before.

\item[d)] We say that $X$ is \textit{locally left filtering} if the connected components are left filtering, i.e. if $X$ satisfies conditions PS 1 and PS 2 (``pseudo-cofiltering'' in the terminology of SGA 4). \textbf{NB} If $X$ is locally left filtering, the same holds for every $X'$ étale over $X$, for this is obvious for PS 1 which concerns the $X/x$, and for PS 2 this follows from the fact that $X' \longrightarrow X$ is faithful. So if $X$ is locally cofiltering, the $X/a$ are filtering, but the converse is false.
\end{enumerate}

\printedpage{11}
% head: A. Grothendieck : Dérivateurs, Ch. II, Édition des Universités de Montpellier II et Paris VII
% foot: Version 8 juillet 2001

\noindent We have the dual notions: \textit{weakly right filtering} (or simply weakly filtering), and \textit{right filtering} or simply filtering. (Dual notions of a), c) above.) For the dual notion of b), namely that the $x\backslash X$ be weakly filtering, we can no longer call it ``locally weakly filtering'', but it would be \textit{co}locally weakly filtering.

\noindent\textbf{Example.} $X = \Bigl(0 \xrightarrow{\gamma} 1 \underset{\beta'}{\overset{\alpha'}{\rightrightarrows}} 2 \underset{\beta}{\overset{\alpha}{\rightrightarrows}} 3\Bigr)$ with $(\alpha\alpha' = \beta\beta') \neq (\alpha\beta' = \beta\alpha')$.

\msp{page 17}

\noindent Only the notion a) (weakly left filtering (\footnote[22]{pseudo-filtering.})) depends only on $\operatorname{Ord} X$ - but of course if $X$ is \textit{pre}ordered, all three depend only on $\operatorname{Ord} X$. Let us note that saying it depends only on $\operatorname{Ord} X$, or only on $\mathrm{Esp}\,\mathrm{sep}X$, amounts to the same thing, since the two determine each other. Condition a) means that $\operatorname{Esp} X$, or also $\operatorname{Top}(X)$, is \textit{irreducible.} We shall say that $X$ is \textit{irreducible} (instead of pseudo-filtering). The other three notions cannot be expressed in terms of $\operatorname{Esp} X$ or $\operatorname{Ord} X$ alone, but they can be expressed in terms of the $\mathrm{Topos}\, X$. Condition b) ($X$ locally weakly filtering) means that $\operatorname{Top}(X)$ is \textit{locally irreducible,} and we shall also say that $X$ is \textit{locally irreducible} (instead of locally weakly filtering). This means that for every $X'$ étale over $X$ and 0-connected, $X'$ is weakly left filtering, i.e. is irreducible.

It remains to express the filtering condition itself in terms of $\operatorname{Top} X$. I recall in this connection:

\noindent\textbf{Proposition 2.3.} \textit{The following conditions on a topos $\sX$ are equivalent (when $S$ (as a ``site'') denotes a generating subset of} $\mathrm{Fais}(\sX) = S^{\wedge}$, consisting of 0-connected objects (\footnote[23]{thus $\sX$ is locally 0-connected.}).\textit{)}
\begin{enumerate}[label={},leftmargin=3em]
\item[\textit{a)}] \textit{If $U$, $V$ are two $0$-connected objects of $\operatorname{Ob} \sX$, $U \times V$ is $0$-connected.}
\item[\textit{b)}] \textit{If $U$ is a $0$-connected object of $\sX^{\wedge}$, then $U \longrightarrow e$ is ``$W_0$-aspheric'', i.e. for every $V$, the projection $V \times U \longrightarrow V$ induces $\pi_0(U \times V) \xrightarrow{\sim} \pi_0(U)$.}
\item[\textit{c)}] \textit{If $U$, $V$} \texttt{[are]} \textit{in $\sX^{\wedge}$, then}
\[
\pi_0(U \times V) \xrightarrow{\sim} \pi_0(U) \times \pi_0(V).
\]
\item[\textit{d)}] \textit{As in a), with $U$, $V$ in $S$. (But we assume the objects of $S$ to be $0$-connected.)}
\end{enumerate}
\noindent(\footnote[24]{The proof should be made explicit together with that of 2.9 below (p.~25).}). % [sic] démonstation -> démonstration, letter r missing

\msp{page 18}

\noindent Moreover, these conditions imply that $\sX$ is connected, hence 0-connected if $\sX \neq \emptyset$. This moreover implies that $\sX$ is $W_{\infty}$-aspheric, cf. Pursuing Stacks §35 (Corollary p.~52) (\footnote[25]{\textbf{False,} loc. cit. says nothing of the kind.}).

It is essentially Proposition 1 of Pursuing Stacks §65, cf. loc. cit. p.~170-171 (\footnote[26]{But there I restrict myself to topoi of the form $\operatorname{Top}(A)$, with $A$ in Cat.}). Cf. also Pursuing Stacks, §35, pages 50-55.

\printedpage{12}
% head: A. Grothendieck : Dérivateurs, Ch. II, Édition des Universités de Montpellier II et Paris VII
% foot: Version 8 juillet 2001

\noindent\textbf{Definition 2.4.} We say that $\sX$ is \textit{totally $0$-connected} if it is locally connected, nonempty and satisfies the equivalent conditions of 2.3. (NB this is a special case of the notion ``totally $W$-aspheric'' of loc. cit. Thus we have the notion of totally $n$-connected.)

Let us note that for an object $X$ of Cat, saying that $\operatorname{Top}(X)$ is totally 0-connected amounts to saying that $X \neq \emptyset$, and that for $x, y \in \operatorname{Ob} X$, $x \times y$ in $X^{\wedge}$ is 0-connected, i.e. $X/x \times y$ is 0-connected (\footnote[27]{This also means that $\mathrm{diag}_X : X \longrightarrow X \times X$ is co-cofinal, i.e. $X^o \longrightarrow X^o \times X^o$ is cofinal, and implies that the same holds for $X \longrightarrow X^E$ for every \textit{finite} set $E$.}). It is a very strong condition. We shall also say that $X$ is \textit{totally $0$-connected.} Let us note here

\noindent\textbf{Proposition 2.5.} \textit{Let $\sX$ be a topos. If $\sX$ is totally $0$-connected, then $\sX$ is irreducible. (In fact the conclusion holds, tautologically, assuming only that the product of two $0$-connected objects of $\sX^{\wedge}$ is nonempty, and it even suffices to assume this for subobjects of the final object $e_{\sX}$.) In particular, if $\sX = \mathit{Top}(X)$, where $X$ is in} Cat\textit{, then $X$ is weakly left filtering.}

But of course, this does \textit{not} imply that $\sX$ is locally irreducible.

\msp{page 19}

\noindent For example the category $\boldsymbol{\Delta}$, or any small category stable under binary products, is totally 0-connected, but $\boldsymbol{\Delta}$ is not locally irreducible (the induced categories $\boldsymbol{\Delta}/\boldsymbol{\Delta}^n$, for $n > 0$, are visibly \textit{not} irreducible). (But let us recall that if $\sX = \operatorname{Top}(T)$, $T$ a topological space, then $T$ irreducible implies that $\sX$ is also \textit{locally} irreducible. So $T$ totally 0-connected implies $\operatorname{Top}(T)$ globally \textit{and} locally irreducible, i.e. for every $T'$ étale over $T$, and 0-connected, $T'$ is irreducible.)

\noindent\textbf{Theorem 2.6.} \textit{Let $\sX$ be a locally connected topos, $S \subseteq \sX^{\wedge}$ a small generating subset consisting of $0$-connected objects. Equivalent conditions (}except \textit{c, c$'$, which seem} (\footnote[28]{\textbf{are.}}) \textit{weaker, cf. proof).)}
\begin{enumerate}[label={},leftmargin=3em]
\item[\textit{a)}] \textit{Every induced topos $\sX/U$ over $U$ $0$-connected in $\sX^{\wedge}$ is totally $0$-connected.}
\item[\textit{a$'$)}] \textit{Every fiber product $U \times_W V$ of $0$-connected objects in $\sX^{\wedge}$ is $0$-connected.}
\item[\textit{b)}] \textit{$\sX$ is locally irreducible, and its connected components are totally $0$-connected.}
\item[\textit{c)}] \textit{As in a), with $U$ in $S$, i.e. the $\sX/U$ ($U \in \operatorname{Ob} S$) are totally $0$-connected.}
\item[\textit{c$'$)}] \textit{As in a$'$), with $U$, $V$, $W$ in $S$.}
\item[\textit{d)}] \textit{$S$ is locally left filtering, i.e. $S$ satisfies the two conditions}

PS$^*$ 1\ \ \textit{Every diagram}
\begin{tikzcd}[ampersand replacement=\&, row sep=tiny, column sep=small, baseline=(\tikzcdmatrixname.center)]
a \arrow[dr] \& \\
 \& x \\
b \arrow[ur] \&
\end{tikzcd}
\textit{in $S$ fits into}
\begin{tikzcd}[ampersand replacement=\&, row sep=tiny, column sep=small, baseline=(\tikzcdmatrixname.center)]
 \& a \arrow[dr] \& \\
c \arrow[ur, dotted] \arrow[dr, dotted] \& \& x \\
 \& b \arrow[ur] \&
\end{tikzcd}
\textit{which is commutative.}

PS$^*$ 2\ \ \textit{Every diagram $a \rightrightarrows b$ fits into a commutative diagram $c \longrightarrow a \rightrightarrows b$.}
\end{enumerate}
\noindent(\footnote[29]{NB We have equivalence of \tikz[baseline=(n.base)]\node[draw, rounded corners, inner sep=1.5pt](n){a), a$'$), b), d),}; which imply c), c$'$), which are equivalent. But it seems (\footnotemark[30]) that c), c$'$) are strictly weaker. Cf. however the variants e), e$'$) in 2.12 (p.~28, 29).}). % [sic] plus faible -> plus faibles, adjective does not agree with the plural subject c), c')
\footnotetext[30]{It turns out. It is indeed true, c) and c$'$) do \textit{not} imply the other conditions, cf. marginal annotation
% footnote 30 continues at the foot of printed page 13:
p.~28 (case of $X = B_G$).}

\printedpage{13}
% head: A. Grothendieck : Dérivateurs, Ch. II, Édition des Universités de Montpellier II et Paris VII
% foot: Version 8 juillet 2001

\msp{page 20}

\noindent These conditions imply that the 0-connected objects of $\sX^{\wedge}$ (in particular the elements of $S$) are $\infty$-connected.

\noindent\textbf{Corollary 2.7} (of the equivalence b $\Longleftrightarrow$ d). \textit{Let $X$ be an object of} Cat. \textit{For $X$ to be left filtering it is necessary and sufficient that $\mathit{Top}(X)$ be totally $0$-connected and locally (and therefore also globally) irreducible (or again that $X$ be totally $0$-connected and that for every $x \in \operatorname{Ob} X$, the localized category $X/x$ be irreducible (it is even totally $0$-connected), or again that for two objects $a$, $b$ in $X$, $a \times b \in X^{\wedge}$ be $0$-connected, and for}
\begin{tikzcd}[ampersand replacement=\&, row sep=tiny, column sep=small, baseline=(\tikzcdmatrixname.center)]
a \arrow[dr] \& \\
 \& x \\
b \arrow[ur] \&
\end{tikzcd}
\textit{in $X$, $a \times_x b$ (in $X^{\wedge}$) be $0$-connected)} (\footnote[31]{Example of $X$ totally 0-connected, but not cofiltering, i.e. not locally irreducible? The category $\boldsymbol{\Delta}$ of standard simplices is totally 0-connected, but is not cofiltering because $\boldsymbol{\Delta}^0 = e \underset{\partial_1}{\overset{\partial_0}{\rightrightarrows}} \boldsymbol{\Delta}^1$ cannot be equalized! So \texttt{[the opposite category]} $\boldsymbol{\Delta}^o = I$ is such that $I \longrightarrow I^E$ is cofinal for every finite set $E$, but $I$ not filtering.}).

\noindent\textsc{Proof} of 2.6. We have the diagram of $\pm$ tautological implications
\[
\begin{tikzcd}[ampersand replacement=\&, column sep=large, row sep=large]
 \& \mathrm{a} \arrow[r, Leftrightarrow] \arrow[dl, Rightarrow, "{\substack{\text{by}\\ \text{2.5}}}"'] \arrow[d, Rightarrow] \& \mathrm{a}' \arrow[d, Rightarrow] \\
\mathrm{b} \arrow[d, Rightarrow] \& \mathrm{c} \arrow[r, Leftrightarrow, "{(\Leftarrow \text{ by 2.3.})}"] \& \mathrm{c}' \arrow[dll, Rightarrow] \\
\text{PS 1.} \& \&
\end{tikzcd}
\]
Let us nevertheless show the two implications c$'$) $\Longrightarrow$ PS 1 and b $\Longrightarrow$ PS 1.

\msp{page 21}

\noindent This is contained in the

\noindent\textbf{Corollary 2.8.} \textit{Let $\sX$, $S$ be as in 2.6. Equivalent conditions:}
\begin{enumerate}[label={},leftmargin=3em]
\item[\textit{a)}] \textit{$\sX$ is locally irreducible, i.e. the induced topoi $\sX/U$, where $U$ is $0$-connected, are irreducible.}
\item[\textit{b)}] \textit{Every fiber product $U \times_W V$ of $0$-connected objects $U$, $V$ is ``nonempty''.}
\item[\textit{b$'$)}] \textit{Every fiber product $U \times_W V$, with $U$, $V$ ``nonempty'', $W$ $0$-connected, is nonempty.}
\item[\textit{c)}] \textit{The $\sX/x$ with $x \in S$ are irreducible.}
\item[\textit{d)}] \textit{As in b), with $U$, $V$, $W$ in $S$.}
\item[\textit{e)}] \textit{$S$ is locally weakly left filtering, i.e. satisfies PS$^*$ 1.}
\end{enumerate}
\noindent Here the $\pm$ obvious implications are

\end{document}
